\subsection{The finite elliptic Euler complexes and the actual kernel}
\label{subsec:cm-elliptic-euler-kernel}

There is a finite-level description of the precise arithmetic
image entering \(J_A\).  Retain the finite-\(\mathcal X_2^+\)
hypothesis and write
\[
 V_n=E[p^n],\qquad
 \alpha_n:A\longrightarrow\mathbb Z_p[V_n](1),\qquad
 A_n=\operatorname{im}\alpha_n.
\]
The transition on \(V_n\) is multiplication by \(p\).
Each \(A_n\) is its actual compact, finite free
\(\mathbb Z_p\)-lattice.

Use Goncharov's elliptic units \(\theta_E(x)\), and let
\(C_1(n)\) be their \(\mathbb Q_p\)-span for
\(0\ne x\in V_n\), inside \(L_n^\times\otimes\mathbb Q_p\).
Here \(B_2(F)\) denotes Goncharov's rational Bloch group.
It is generated by symbols \(\{z\}\), \(z\in\mathbb P^1(F)\),
with differential \(\delta\{z\}=(1-z)\wedge z\), set equal to
zero at \(0,1,\infty\). Its relations identify the two
specializations of \(\sum_i\{f_i\}\) along a connected curve
whenever \(\sum_i(1-f_i)\wedge f_i=0\) in its function field;
see \cite[Section 2.1, p.~10]{GoncharovEuler}.
For number fields this agrees rationally with the five-term
presentation through the canonical map stated there.
Let \(C_2(n)\) be the \(\mathbb Q_p\)-span of his
double elliptic units \(\theta_E(a,b,c)\), with
\(a,b,c\in V_n\) and \(a+b+c=0\), in
\(B_2(L_n)\otimes\mathbb Q_p\).
The ray field and these theta-value spans are described in
\cite[Definition 5.2, Proposition 5.3 and Definition 5.4,
p.~27]{GoncharovEuler}.
Its differential is
\begin{equation}\label{eq:cm-elliptic-triangle-delta}
 \delta\theta_E(a,b,c)
   =\theta_E(a)\wedge\theta_E(b)
    +\theta_E(b)\wedge\theta_E(c)
    +\theta_E(c)\wedge\theta_E(a),
\end{equation}
by \cite[Theorem 2.10(a), equation (38), p.~15]{GoncharovEuler}.
The differential of a triple with a zero entry vanishes;
no value of \(\theta_E(0)\) is needed in its image.
Put
\[
 \mathcal R_n^{\rm ell}
   =\operatorname{coker}\bigl(C_2(n)
                          \xrightarrow{\delta}\bigwedge^2C_1(n)\bigr),
 \qquad
 \mathcal D_n=\ker\bigl(J_n:\bigwedge^2A_n
                   \longrightarrow\mathbb Z_p[V_n\times V_n](2)\bigr),
\]
where \(J_n\) is the same pair of finite additive shears
as in Theorem~\ref{thm:cm-measure-class-two}.

\begin{theorem}\label{thm:cm-elliptic-euler-kernel}
There are natural equivariant identifications
\[
 A_n[1/p]=\operatorname{Hom}_{\mathbb Q_p}
                         (C_1(n),\mathbb Q_p(1)),
 \qquad
 \mathcal D_n[1/p]=\operatorname{Hom}_{\mathbb Q_p}
                         (\mathcal R_n^{\rm ell},\mathbb Q_p(2)).
\]
The transition \(\mathcal D_{n+1}\to\mathcal D_n\)
is induced by the actual pushforward
\(A_{n+1}\to A_n\).  Rationally, it is dual to the
natural map \(\mathcal R_n^{\rm ell}\to
\mathcal R_{n+1}^{\rm ell}\) induced by field extension
and the elliptic-unit distributions.  Finally,
\begin{equation}\label{eq:cm-J-kernel-elliptic-limit}
                       \ker J_A=\varprojlim_n\mathcal D_n.
\end{equation}
Here the inverse limit retains the specified integral
lattices \(\mathcal D_n\).
\end{theorem}

\begin{proof}
We first check the normalization linking the two primary
sources.  Scale the period lattice to
\(L=\mathbb Z\tau+\mathbb Z\), and write
\(z=a\tau+b\), \(q=e^{2\pi i\tau}\),
\(\xi=e^{2\pi iz}\).  Put
\(\eta_1=\eta(\tau)\), \(\eta_2=\eta(1)\).
The sigma product and the Legendre relation give
\[
 \sigma(z,L)=\frac{e^{\eta_2z^2/2}}{2\pi i}
       \xi^{-1/2}(\xi-1)
       \prod_{m\ge1}
       \frac{(1-q^m\xi)(1-q^m\xi^{-1})}{(1-q^m)^2},
 \qquad \eta_2\tau-\eta_1=2\pi i.
\]
The product identity follows from the zeros and the two
quasiperiods of \(\sigma\), with \(\sigma'(0)=1\)
fixing its constant.  Substitution in Nakamura's
\(\Theta(z,L)=\Delta(L)e^{-6\eta(z)z}\sigma(z,L)^{12}\),
using \(\eta(z)=a\eta_1+b\eta_2\), yields
\[
 \Theta(z,L)=q^{6a^2-6a+1}e^{12\pi ib(a-1)}
    (1-\xi)^{12}
    \prod_{m\ge1}(1-q^m\xi)^{12}(1-q^m\xi^{-1})^{12}.
\]
This is precisely the twelfth power of
\cite[equation (19), p.~7]{GoncharovEuler}.
Homogeneity removes the scaling of the lattice.
The root-of-unity ambiguity from torsion representatives
has zero Kummer character on \(H\).  Thus
\(\kappa_{\Theta(x)}=12\kappa_{\theta_E(x)}\);
the factor \(12\) is invertible because \(p\ge5\).

Nakamura's formulas are stronger than a moment identity.
By \cite[equation (3.5.1), p.~209, and
equations (4.11.3)--(4.11.5), p.~220]{Nakamura1995},
the actual finite coefficients satisfy
\begin{equation}\label{eq:cm-alpha-theta-finite}
 \alpha_n(h)(x)=12\kappa_{\theta_E(x)}(h)\quad(x\ne0),
 \qquad
 \alpha_n(h)(0)=-n\kappa_{p^{12}}(h).
\end{equation}
The latter formula also uses
\cite[Proposition (3.8), p.~210]{Nakamura1995}.
Ishii's special-lift measure is this measure, with basis
compatibility recorded in
\cite[Remark 2.6.1(2), p.~8]{Ishii}.
CM-unit invariance and the ray-field description in
\cite[p.~27]{GoncharovEuler} place these rational theta
classes over \(L_n\).

We verify that their restriction to \(H\) loses no
nonzero rational class.  Set
\(G_n=\operatorname{Gal}(\Omega/L_n)\) and
\(\Gamma_n=\operatorname{Gal}(L_\infty/L_n)\).
Finite-field Kummer theory injects the \(p\)-unit
space of \(L_n\) into \(H^1(G_n,\mathbb Q_p(1))\):
the relevant Kummer extensions are pro-\(p\) and
unramified outside \(p\), and hence lie in \(\Omega\).
The restriction kernel to \(H\) is the image of
\(H^1(\Gamma_n,\mathbb Q_p(1))\), which is zero.
Indeed, choose \(\gamma_0\) with
\(\chi_{\rm cyc}(\gamma_0)\ne1\).  Since \(\Gamma_n\)
is abelian, every cocycle \(f\) satisfies
\[
 (\chi_{\rm cyc}(\gamma_0)-1)f(\gamma)
       =(\chi_{\rm cyc}(\gamma)-1)f(\gamma_0),
\]
so it is a coboundary over \(\mathbb Q_p\).
This proves the required Kummer injectivity, including
the possible infinitely \(p\)-divisible kernel in the
base tower.

The nonzero coordinate functionals in
\eqref{eq:cm-alpha-theta-finite} span precisely
\(C_1(n)\); the origin coordinate is their negative
sum, since \(\alpha\) has total mass zero.
A proper rational image of \(A_n[1/p]\) in the dual
of \(C_1(n)\) would have a nonzero annihilator there,
contradicting the Kummer injectivity just proved.
This proves the first identification, while retaining
\(A_n\) as its actual integral lattice.

For \(\eta\in\bigwedge^2A_n[1/p]\), let
\(B_\eta\) denote the alternating pairing on \(C_1(n)\)
under that identification.  The inverse-image formula
for the shears is
\[
 (J_n\eta)(x,y)=\eta(x,y-x)+\eta(x-y,y)-\eta(x,y).
\]
When \(x,y,x-y\ne0\), evenness
\(\theta_E(-u)=\theta_E(u)\) and
\eqref{eq:cm-alpha-theta-finite} give
\begin{equation}\label{eq:cm-J-elliptic-triangle}
 (J_n\eta)(x,y)
       =144 B_\eta\bigl(\delta\theta_E(x,y-x,-y)\bigr).
\end{equation}
When \(x=0\), \(y=0\), or \(x=y\), the coefficient
vanishes by alternation and evenness, independently
of the origin coefficient.  Every nondegenerate
zero-sum triple occurs in
\eqref{eq:cm-J-elliptic-triangle}.
Thus the rational kernel is exactly the dual of
\(\mathcal R_n^{\rm ell}\), with the second Tate
twist.  Its integral kernel \(\mathcal D_n\) is
closed and \(p\)-saturated because its target is
finite free over \(\mathbb Z_p\).

Here are the transition maps, including the origin.
For \(x\ne0\), the isogeny distribution gives, in
additive notation for units tensored with \(\mathbb Q_p\),
\[
 \theta_E(x)=\sum_{[p]y=x}\theta_E(y).
\]
This is the ordinary inclusion of the theta class at
the same torsion point into the next field.  Hence
pushforward \(A_{n+1}\to A_n\) is dual to
\(C_1(n)\hookrightarrow C_1(n+1)\).
At zero the exact calculation is
\[
 \sum_{[p]y=0}\alpha_{n+1}(h)(y)
   =-(n+1)\kappa_{p^{12}}(h)
         +\sum_{0\ne y\in E[p]}\kappa_{\Theta(y)}(h)
   =-n\kappa_{p^{12}}(h),
\]
using \cite[Proposition (2.2)(3), p.~204]{Nakamura1995}.
For double units, independence of the torsion exponent
and the sum over preimage triples are
\cite[Lemma 2.12, p.~17, and equation (41), p.~16]{GoncharovEuler}.
Their differential is the exterior square of the
unit inclusion: two chosen preimages determine the
third by the zero-sum condition.  This proves the
commuting diagram of complexes and the asserted
natural map on cokernels.

Finally, compactness and injectivity of \(\alpha\)
identify \(A\) with \(\varprojlim A_n\).
The coefficient quotients modulo \(p^r\), for all
\(n,r\), form a cofinal system of neighborhoods.
Completed tensor and exterior squares therefore give
\(\widehat{\bigwedge}^2A=\varprojlim\bigwedge^2A_n\).
The shear maps commute with these projections,
so taking their kernels proves
\eqref{eq:cm-J-kernel-elliptic-limit}.
\end{proof}

The transition just used has an explicit comparison
with arithmetic norms.  On a torsion point whose two
Tate coordinates are units,
\(\operatorname{Gal}(L_{n+1}/L_n)\) acts by
\((1+p^nu,1+p^nv)\), \(u,v\in\mathbb F_p\), and
its orbit is the entire \(p^2\)-element \([p]\)-fiber.
The distribution product is then the field norm.
If one coordinate is divisible by \(p\) and the other
is a unit, the arithmetic orbit has \(p\) elements,
each repeated \(p\) times in the norm, whereas the
isogeny fiber has \(p^2\) elements.  The all-chart
transition in the theorem is the isogeny distribution,
with the displayed correction at zero.

For the Gaussian and Eisenstein fields,
\cite[Theorem 5.16, p.~34]{GoncharovEuler} supplies a
surjective morphism from the Bianchi modular complex
at the ideal \((p^n)\) to the complex with target
\(\bigwedge^2\widehat C_1(n)\), where
\(\widehat C_1(n)=C_1(n)\oplus\mathbb Q_p\theta_E(0)\).
Compose its degree-two map with the exterior square
of the projection \(\widehat C_1(n)\to C_1(n)\)
that sends the formal \(\theta_E(0)\) to zero.
The differentials of degenerate triangles vanish,
so this is a surjective morphism to the exact Euler
complex used above.  One can equally send the formal
origin to its actual normalized class \(-n[p]\).
At a prime ideal \(\mathfrak P\),
\cite[Theorem 1.2, p.~9]{GoncharovEuler} identifies
\[
 \operatorname{coker}\bigl(C_2(\mathfrak P)
       \longrightarrow\bigwedge^2C_1^{\rm un}(\mathfrak P)\bigr)
       =H^2_{\rm cusp}(\Gamma_1(\mathfrak P),\mathbb Q),
\]
where \(C_1^{\rm un}\) is the global-unit space.
Since the differential lands in that smaller exterior
square, this gives a canonical subspace of the full
unit cokernel after extension of scalars.  The
global-unit qualification is part of this precise map.

The remaining persistence issue has a concrete form.
Injectivity of the natural map on cokernels requires
\[
 \delta C_2(n+1)\cap\bigwedge^2C_1(n)=\delta C_2(n).
\]
The inclusions of unit spaces alone do not establish
this equality.  Moreover,
\eqref{eq:cm-J-kernel-elliptic-limit} retains the
integral lattices, so rational finite-level classes
must have compatible bounded duals to contribute to
the actual kernel.  The theorem supplies this exact
finite-level and inverse-limit comparison without
inferring such persistence from a dimension count.
